% =============================================================================
% Chapter 1: Introduction
% Hebrew primary, English technical terms in \en{}
% XeLaTeX : requires \selectlanguage{hebrew} at top
% =============================================================================

\selectlanguage{hebrew}

\section{מבוא: ניווט רחפנים בסביבות מוגבלות ראות}
\label{sec:ch01_intro}

\subsection{רקע ומוטיבציה: האתגר של ניווט אוטונומי בתנאי ראות קשים}
\label{sec:ch01_motivation}

רחפנים אוטונומיים (\en{Unmanned Aerial Vehicles, UAVs}) הפכו לכלי מרכזי במגוון רחב של יישומים אזרחיים וצבאיים, החל ממיפוי חקלאי וצילום אווירי ועד לחיפוש והצלה במבנים שקרסו. עם זאת, היכולת של רחפנים לנווט באופן אוטונומי בסביבות מוגבלות ראות נותרה אתגר משמעותי שלא נפתר במלואו. סביבות כגון מנהרות מלאות עשן, מחסנים חשוכים, יערות עם עלווה צפופה, ומבנים שקרסו לאחר רעידת אדמה מציבות קשיים ייחודיים למערכות ניווט קונבנציונליות. \cite{Thrun2005ProbRobotics}

מערכות ניווט מסורתיות מסתמכות על חיישנים אופטיים כגון מצלמות ו\en{LiDAR} (\en{Light Detection and Ranging}) למיפוי הסביבה ומיקום הרחפן בתוכה. חיישנים אלה רגישים מאוד לתנאי הראות: מצלמות דורשות תאורה מספקת וניגודיות חזותית, בעוד ש\en{LiDAR} מסתמך על החזרת אור לייזר ממשטחים, תהליך המופרע באופן משמעותי על ידי עשן, אבק, או ערפל. במנהרה מלאת עשן, לדוגמה, טווח ה\en{LiDAR} יכול לרדת מ-50 מטר לפחות מ-2 מטר, מה שהופך את הניווט לבלתי אפשרי. \cite{Thrun2005ProbRobotics}

המוטיבציה המרכזית של מחקר זה היא לפתח מערכת ניווט לרחפנים המסוגלת לפעול באופן אמין בכל תנאי הראות, תוך השראה ביולוגית ממערכת האקולוקציה של עטלפים. עטלפים, היונקים היחידים המסוגלים לעוף, פיתחו במשך מיליוני שנים מערכת ניווט אקוסטית מתוחכמת המאפשרת להם לנווט בצורה מדויקת בחושך מוחלט, ביערות צפופים, ואף במערות מלאות עשן. מערכת זו, המכונה אקולוקציה (\en{echolocation}), מבוססת על פליטת פולסים על-קוליים וניתוח ההדים החוזרים מהסביבה. \cite{GriffinBatEcholocation} \cite{Simmons1979BatSonar}

\subsection{מערכת האקולוקציה העטלפית: השראה ביולוגית}
\label{sec:ch01_bio_inspiration}

עטלפים משתמשים במערכת אקולוקציה מתוחכמת הכוללת מספר מנגנונים משלימים. ראשית, הם פולטים פולסים על-קוליים בתדרים הנעים בין 20 קילו-הרץ ל-120 קילו-הרץ, תלוי במין ובסביבה. עטלפי פרסה (\en{Rhinolophidae}) משתמשים בפולסים בתדר קבוע (\en{Constant Frequency, CF}) עם אפנון תדר (\en{Frequency Modulation, FM}) בסוף הפולס, בעוד שעטלפים אחרים משתמשים בפולסי \en{FM} בלבד. הפולסים מופקים על ידי מיתרי הקול ומוקרנים דרך הפה או האף, תלוי במין. \cite{Schnitzler1968DSC}

שנית, עטלפים מבצעים סריקת ראש (\en{head scanning}) תוך כדי תעופה, תוך שינוי כיוון הראש בזווית של עד 30 מעלות משני צידי קו הטיסה. סריקה זו מאפשרת להם לקבל מידע אזימוטלי על הסביבה גם עם מערך קולט דו-אוזני בלבד. תדר סריקת הראש נע בין 5 הרץ ל-15 הרץ, תלוי במהירות הטיסה ובמורכבות הסביבה. \cite{MossEcholocation}

שלישית, עטלפים מבצעים פיצוי תדר דופלר (\en{Doppler Shift Compensation, DSC}) אדפטיבי. כאשר עטלף טס לעבר מטרה, תדר ההד החוזר גבוה מתדר הפולס המקורי בשל אפקט דופלר. עטלפי פרסה מתאימים את תדר הפולס כך שתדר ההד החוזר יישאר בתחום הרגישות המרבית של מערכת השמיעה שלהם, שהיא צרה מאוד (כ-100 הרץ). מנגנון זה מאפשר להם לזהות תנועה איטית של טרף (כגון חרקים) על רקע צמחייה נייחת. \cite{Schuller1974DSC}

רביעית, כנפי העטלף משפיעות על תבנית הקרינה האקוסטית. מחקרים הראו כי תנופת הכנפיים גורמת לשינוי בתדר הדופלר הנקלט, וכי עטלפים משתמשים במידע זה להערכת מהירות הטיסה שלהם ביחס לסביבה. מודל זה, המכונה \en{wing-beat modulation}, מספק מידע נוסף על מצב התעופה שאינו זמין מחיישנים אינרציאליים בלבד. \cite{GriffithBatEcholocation}

\subsection{סקירת ספרות: מערכות ניווט קיימות לרחפנים}
\label{sec:ch01_literature}

הספרות המחקרית בתחום ניווט רחפנים בסביבות מוגבלות ראות כוללת מספר גישות מרכזיות. הגישה הראשונה מבוססת על מיזוג חיישנים אינרציאליים (\en{Inertial Measurement Unit, IMU}) עם חיישנים אקוסטיים. מחקרים מוקדמים בתחום זה הראו כי שימוש בחיישן על-קולי יחיד יכול לספק מידע טווח בסביבות פנים-מבנה, אך הרזולוציה האזימוטלית מוגבלת ביותר. \cite{Thrun2005ProbRobotics}

הגישה השנייה מבוססת על \en{SLAM} (\en{Simultaneous Localization and Mapping}) אקוסטי, הבונה מפה של הסביבה תוך כדי מיקום הרחפן בתוכה. \en{ORB-SLAM3}, המערכת המתקדמת ביותר בתחום ה\en{SLAM} החזותי, משיגה דיוק מיקום של 0.048 מטר בתנאי ראות טובים, אך נכשלת לחלוטין בעשן ובחושך. \en{FAST-LIO2}, מערכת \en{SLAM} מבוססת \en{LiDAR}, משיגה דיוק של 0.12 מטר בתנאי ראות טובים, אך ביצועיה יורדים משמעותית בעשן. \cite{MurArtal2015ORB} \cite{Zhang2014LOAM}

הגישה השלישית מבוססת על למידה עמוקה לעיבוד הדים אקוסטיים. רשתות קונבולוציה (\en{CNN}) ורשתות נוירונים גרף (\en{GNN}) הוכחו כיעילות בסיווג סוגי מכשולים מזיהוי הדים על-קוליים. עם זאת, רוב המחקרים בתחום זה התמקדו בזיהוי מכשולים ולא בניווט מלא הכולל מיפוי ומיקום. \cite{CrewAIDocs}

\subsection{ניסוח פורמלי של בעיית הניווט}
\label{sec:ch01_formal_problem}

בעיית הניווט האוטונומי בסביבות מוגבלות ראות מנוסחת באופן פורמלי כבעיית אופטימיזציה סטוכסטית. בהינתן מצב התחלתי $\mathbf{x}_0$, מפה לא ידועה $\mathcal{M}$, ויעד $\mathbf{x}_g$, על הרחפן למצוא מסלול $\tau = \{\mathbf{x}_0, \mathbf{x}_1, ..., \mathbf{x}_T\}$ הממזער פונקציית עלות תוך עמידה באילוצי בטיחות. פונקציית העלות הכוללת מוגדרת כ:

\begin{equation}
J(\tau) = \mathbb{E}\left[ \sum_{k=0}^{T-1} $c(\mathbf{x}_k, \mathbf{u}_k)$ + $c_$f(\mathbf{x}_T)$$ \right]
\label{eq:ch01_cost_function}
\end{equation}

כאשר $c(\mathbf{x}_k, \mathbf{u}_k)$ היא עלות הצעד בזמן $k$ (כוללת מרחק ליעד, סיכון התנגשות, וצריכת אנרגיה), $c_$f(\mathbf{x}_T)$$ היא עלות סופית (מרחק ליעד הסופי), והתוחלת היא על פני אי-הוודאות במצב הרחפן ובמפת הסביבה. אילוצי הבטיחות מבטיחים כי המרחק לכל מכשול $\mathbf{m}_i$ במפה יהיה גדול ממרחק בטוח מינימלי $d_{\text{safe}}$:

\begin{equation}
\| \mathbf{x}_k - \mathbf{m}_i \| \geq d_{\text{safe}} \quad \forall k, \forall i
\label{eq:ch01_safety_constraint}
\end{equation}

בעיה זו היא NP-קשה באופן כללי, ולכן נדרשת גישה היוריסטית מבוססת תכנון מסלול דגימה (\en{sampling-based planning}) כגון \en{RRT*}. \cite{Thrun2005ProbRobotics}

\subsection{תרומות המאמר}
\label{sec:ch01_contributions}

מאמר זה מציג ארבע תרומות מרכזיות:

\begin{enumerate}
    \item מודל ביומימטי שלם של מערכת האקולוקציה העטלפית, הכולל פיצוי תדר דופלר אדפטיבי, סריקת ראש, ומודל קרינת כנפיים. המודל מתרגם עקרונות ביולוגיים לאלגוריתמים חישוביים הניתנים להטמעה על פלטפורמה מוגבלת משאבים.
    \item שיטת מיזוג חיישנים מבוססת \en{EKF} (\en{Extended Kalman Filter}) עם אומדן קווריאנס אדפטיבי, המותאמת לחיישן על-קולי, \en{IMU}, וחיישן זרימה אופטית. השיטה מזהה השפלה של חיישן באמצעות מבחן כי-ריבוע ומתאימה את משקל החיישן באומדן המצב בהתאם.
    \item אלגוריתם \en{SLAM} אקוסטי הבונה מפת רשת תאים הסתברותית תלת-ממדית תוך שימוש ב\en{GNN} להתאמת הדים. האלגוריתם מאפשר מיפוי ומיקום בו-זמניים בסביבות מוגבלות ראות.
    \item מתכנן מסלול מבוסס \en{RRT*} (\en{Rapidly-exploring Random Tree Star}) המותאם לאי-ודאות חישה משתנה, המשלב פונקציית עלות רב-מודאלית הכוללת בטיחות, יעילות אנרגטית, ואיכות חישה. \cite{Thrun2005ProbRobotics} \cite{CrewAIDocs}
\end{enumerate}

\subsection{מבנה המאמר}
\label{sec:ch01_structure}

המאמר מאורגן כדלקמן. פרק 2 מציג את הבסיס הביולוגי של מערכת האקולוקציה העטלפית ואת המודל הביומימטי המוצע. פרק 3 מתאר את מודל החיישן העל-קולי ואת עיבוד האותות האקוסטי. פרק 4 מציג את אלגוריתם ה\en{SLAM} האקוסטי. פרק 5 מתאר את שיטת מיזוג החיישנים הרב-מודאלי. פרק 6 מציג את האלגוריתם המלא לניווט אוטונומי. פרק 7 מתאר את ארכיטקטורת המערכת המשולבת וההטמעה בזמן אמת. פרק 8 מציג תוצאות סימולציה וניסויים. פרק 9 מסכם את המאמר, דן במגבלות, ומציע כיווני מחקר עתידיים.

\subsection{סימונים ומשתנים}
\label{sec:ch01_notation}

לאורך המאמר, נשתמש בסימונים הבאים. וקטורים מסומנים באותיות מודגשות קטנות ($\mathbf{x}$), מטריצות באותיות מודגשות גדולות ($\mathbf{A}$), וסקלרים באותיות רגילות ($x$). הזמן מסומן ב-$k$ עבור צעד בדיד. מצב הרחפן בזמן $k$ מיוצג על ידי הווקטור $\mathbf{x}_k = [\mathbf{p}_k, \mathbf{q}_k, \mathbf{v}_k, \mathbf{b}_k]^T$, כאשר $\mathbf{p}_k$ הוא מיקום תלת-ממדי, $\mathbf{q}_k$ הוא קווטרניון ייחוס, $\mathbf{v}_k$ היא מהירות, ו-$\mathbf{b}_k$ הוא וקטור הטיות החיישנים. \cite{Thrun2005ProbRobotics}

% =============================================================================
% End of Chapter 1
% =============================================================================